\documentclass[10pt]{article}
\usepackage[margin=1in]{geometry}
\usepackage{amsmath,amssymb,amsthm,hyperref}
\newtheorem{theorem}{Theorem}
\title{A two-dimensional operator of minimal isometry order three}
\author{ziangni-sys}
\date{Conjecture 00000008576}
\begin{document}
\raggedright
\maketitle
\noindent Both source versions assert that the minimal order $m$ of an
$m$-isometric operator is bounded above by the dimension of its space.
We disprove this clause with a genuine operator on the complex Euclidean
Hilbert space $H=\mathbb C^2$.

\section*{Definitions and the actual operator}
The inner product and norm on $H$ are
\[
 \langle v,w\rangle=\overline v_0w_0+\overline v_1w_1,
 \qquad \|v\|^2=|v_0|^2+|v_1|^2.
\]
The standard basis is orthonormal. For a bounded linear operator $T$, define
\[
 B_m(T)=\sum_{j=0}^m(-1)^{m-j}\binom mj(T^*)^jT^j.
\]
The usual $m$-isometry condition is $B_m(T)=0$, with $m\geq1$.
Equivalently, the corresponding alternating binomial sum of
$\|T^jv\|^2$ vanishes for every $v$. The adjoint form follows from
$\langle v,(T^*)^jT^jv\rangle=\|T^jv\|^2$ and the polarization identity.
The alternative sign convention $(-1)^j$ multiplies the whole identity by
$(-1)^m$ and gives the same zero condition.

Let $T(v_0,v_1)=(v_0+v_1,v_1)$. Its actual matrix is
\[
 J=\begin{pmatrix}1&1\\0&1\end{pmatrix}.
\]
Every linear map on this finite-dimensional Hilbert space is bounded.
Induction using matrix multiplication gives
\[
 J^j=\begin{pmatrix}1&j\\0&1\end{pmatrix},\qquad
 (J^j)^*J^j=\begin{pmatrix}1&j\\j&j^2+1\end{pmatrix}
 \quad(j\geq0).
\]
Here conjugate transpose is precisely the Hilbert adjoint in the standard
orthonormal basis.

\begin{theorem}
$T$ is a $3$-isometry, is neither a $1$- nor a $2$-isometry, and therefore
has minimal positive isometry order $3>\dim_{\mathbb C}H=2$.
\end{theorem}
\begin{proof}
Write $G_j=(J^j)^*J^j$. Actual matrix computations yield
\[
 B_1(J)=G_1-G_0=\begin{pmatrix}0&1\\1&1\end{pmatrix}\ne0,
 \qquad
 B_2(J)=G_2-2G_1+G_0=\begin{pmatrix}0&0\\0&2\end{pmatrix}\ne0.
\]
For order three,
\[
 B_3(J)=G_3-3G_2+3G_1-G_0=0:
\]
each entry of $G_j$ is a polynomial in $j$ of degree at most two, whose
third alternating binomial difference vanishes. These are identities of
full matrices, so they hold for all vectors, not just a sample vector.
The only positive orders below three are one and two, both excluded above.
Thus the minimal positive order is three, contradicting the asserted
upper bound by dimension two.
\end{proof}

\newpage
\section*{Formal Hilbert-space correspondence}
The Lean project uses the actual Mathlib space
\texttt{EuclideanSpace} $\mathbb C$ \texttt{(Fin 2)}, with its Hermitian inner
product and Euclidean norm. It does not use the default supremum norm on a
function space to define an isometry.
\texttt{ambient\_dimension} proves its actual complex \texttt{Module.finrank}
is $2$. The theorem \texttt{jordan\_continuous} also verifies the
continuity of the actual finite-dimensional operator.

The definition \texttt{J} is the actual complex matrix displayed above.
\texttt{jordan\_pow} proves its matrix power formula for every natural $j$.
\texttt{gram m A} is exactly the full finite sum defining $B_m(A)$,
with Mathlib's actual conjugate transpose, matrix powers and multiplication,
and natural binomial coefficients cast into $\mathbb C$.
The theorems \texttt{gram\_one}, \texttt{gram\_two} and
\texttt{gram\_three} verify the three full matrix identities.
Their nonzero witnesses are respectively entry $(0,1)$ and entry $(1,1)$;
there are no asserted spectral or norm values.

Crucially, the project also transports these matrix computations to
actual Hilbert-space linear operators. The definition
\texttt{representation} is the inverse of Mathlib's star-algebra equivalence
\texttt{LinearMap.toMatrixOrthonormal}, taken in the actual standard
orthonormal basis \texttt{EuclideanSpace.basisFun}. This equivalence
preserves scalar multiplication, sums, products, powers and adjoints.
In the linear endomorphism algebra, \texttt{star} is the Hilbert adjoint
and multiplication is composition.

\texttt{operatorGram} is the same binomial identity in that actual
endomorphism algebra. \texttt{representation\_gram} proves that it is the
image of the matrix Gram sum, and \texttt{gram\_operator\_iff} proves
that one vanishes exactly when the other does.
\texttt{jordan\_is\_three} therefore proves that the actual operator is a
$3$-isometry. \texttt{no\_smaller\_positive\_order} quantifies every positive
$m<3$ and excludes that order. The final theorem
\texttt{dimension\_bound\_fails} combines genuine $3$-isometry,
positive-order minimality and the strict actual dimension inequality.

\section*{Scope}
This is a disproof of the explicitly stated minimal-order dimension-bound
conjunct in both languages. The example has positive finite dimension and
positive minimal order. No claim about the separate spectrum, unitary
criterion, or extension-decomposition clauses is needed to disprove the
conjunction. The convention that an $m$-isometry may also satisfy identities
at higher orders causes no ambiguity: we compute its least positive order.

\section*{Reproduction}
The project pins Lean 4.19.0 and Mathlib commit
\begin{center}\small\texttt{c44e0c8ee63ca166450922a373c7409c5d26b00b}.\end{center}
From \texttt{lean/}, run \texttt{lake update}, \texttt{lake exe cache get},
and \texttt{lake build}; then check the entire source with
\begin{center}\small\texttt{lake env lean -DwarningAsError=true Main.lean}.\end{center}
The source prints dependency audits for the main matrix, representation and
minimality theorems. No statements are admitted, and only standard logical
axioms are used.
\end{document}
